\section{Приложения}

\subsection{Краткая сводка отличий от популярных наборов правил}

\noindent\begin{tabularx}{\linewidth}{L{0.4\linewidth}L{0.15\linewidth}L{0.15\linewidth}L{0.15\linewidth}}
	\caption{Отличия правил} \\
	\toprule
	\textbf{Правила} & \textbf{EMA} & \textbf{WRC} & \textbf{RRC} \\
	\endfirsthead
	\toprule
	\textbf{Правила} & \textbf{EMA} & \textbf{WRC} & \textbf{RRC} \\
	\midrule
	\endhead
	\multicolumn{4}{r}{\footnotesize(Продолжение на следующей странице)}
	\endfoot
	\bottomrule
	\endlastfoot
	\multicolumn{4}{c}{\cellcolor{gray!25}Частично отличающиеся правила} \\
	
	
	Риичи, оставшиеся на столе после завершения игры &
	уходят победителю &
	теряются &
	уходят победителю \\
	\midrule
	В случае равенства очков победителей риичи на кону &
	делятся поровну &
	- &
	делятся поровну \\
	\midrule
	Акадоры &
	нет &
	нет &
	есть\\
		\midrule
	Приоритет объявлений &
	рон > пон/кан > чи &
	\textit{рон > пон/кан/чи} &
	рон > пон/кан > чи \\
	\multicolumn{4}{L{0.93\linewidth}}{\small{\textit{Примечание: в WRC в ином случае приоритет определяется по первенству объявления}}} \\
	\midrule
	Атамаханэ &
	нет &
	есть &
	нет\\
	\midrule
	Кто получает выплаты за хонбу в случае дабл-/триплрона &
	все &
	- &
	все (если применимо) \\
	\midrule
	Кто получает палочки риичи на кону в случае дабл-/триплрона &
	первый по ходу игры &
	- &
	первый по ходу игры (если применимо) \\
	\midrule
	Добавление хонбы при абортивной ничьей &
	- &
	- &
	да (если применимо) \\
	\midrule
	Пао на сууканцу &
	нет &
	да &
	нет \\
	\midrule
	Кириаге манган &
	есть&
	есть &
	нет \\
	\multicolumn{4}{L{0.93\linewidth}}{\small{\textit{Пояснение: Кириаге манган --- округление рук 4/30 и 3/60 до мангана}}} \\
	\midrule
	Счетный якуман (казое) &
	нет &
	есть &
	есть \\
	\midrule
	Сложение якуманов &
	нет &
	есть &
	есть \\
	\midrule
	Ограбление закрытого кана при выигрыше в 13 сирот &
	можно &
	нельзя &
	можно \\
	\midrule
	Фу за ветер места и раунда &
	2 фу &
	2 фу &
	4 фу \\
	\midrule
	Объявление риичи без возможности следующего взятия из стены &
	допустимо&
	допустимо&
	недопустимо\\
	\multicolumn{4}{c}{\cellcolor{gray!25}Полностью совпадающие правила} \\
	Стартовые очки &
	\multicolumn{3}{c}{30000} \\
	\midrule
	Ума &
	\multicolumn{3}{c}{15000/5000} \\
	\midrule
	Куитан (открытое тан-яо) &
	\multicolumn{3}{c}{есть} \\
	\midrule
	Атодзуке &
	\multicolumn{3}{c}{есть} \\
	\multicolumn{4}{L{0.93\linewidth}}{\small{\textit{Пояснение: атодзуке --- возможность победы, даже если одно из ожиданий не дает яку.}}} \\
	\midrule
	Куикаэ &
	\multicolumn{3}{c}{недопустимо} \\
	\multicolumn{4}{L{0.93\linewidth}}{\small{\textit{Пояснение: куикаэ --- возможность сбросить тайл, завершающий только что объявленный сет.}}} \\
	\midrule
Абортивные ничьи &
	\multicolumn{3}{c}{нет} \\
	\midrule
	Ока &
	\multicolumn{3}{c}{0} \\
	\midrule
	В случае равенства очков ума: &
	\multicolumn{3}{c}{делится поровну} \\
	\midrule
	После окончания турнирного таймера &
	\multicolumn{3}{c}{доиграть текущую раздачу и сыграть еще одну} \\
	\midrule
	Ренчан назначается в случае если &
	\multicolumn{3}{c}{дилер победил или темпай} \\
	\midrule
	Выплата за хонбу &
	\multicolumn{3}{c}{300} \\
	\midrule
	Суммарные выплаты в случае ничьей &
	\multicolumn{3}{c}{3000} \\
	\midrule
	Дора &
	\multicolumn{3}{c}{есть} \\
	\midrule
	Урадора &
	\multicolumn{3}{c}{есть} \\
	\midrule
	Кандора &
	\multicolumn{3}{c}{есть} \\
	\midrule
	Кан-урадора &
	\multicolumn{3}{c}{есть} \\
	\midrule
	Ренхо &
		\multicolumn{3}{c}{манган} \\
	\midrule
	Выход из временного фуритена &
	\multicolumn{3}{c}{дискард} \\
	\midrule
	Рянхансибари &
	\multicolumn{3}{c}{нет} \\
	\midrule
	Открытие индикатора кандоры при объявлении открытого кана &
	\multicolumn{3}{c}{сразу} \\
	\midrule
	Открытие индикатора кандоры при объявлении закрытого кана &
	\multicolumn{3}{c}{сразу} \\
	\midrule
	Пао на дайсанген &
	\multicolumn{3}{c}{да} \\
	\midrule
	Пао на дайсууши &
	\multicolumn{3}{c}{да} \\
	\midrule
	Секинин барай (ответственность за риншан) &
	\multicolumn{3}{c}{нет} \\
	\midrule
	Выплата хонбы при пао &
	\multicolumn{3}{c}{только набросивший в рон} \\
	\midrule
	Агарияме&
	\multicolumn{3}{c}{нет} \\
	\midrule
	Нагаши манган &
	\multicolumn{3}{c}{нет} \\
		\midrule
	Западные раунды &
	\multicolumn{3}{c}{никогда} \\
	\midrule
	Банкротство &
	\multicolumn{3}{c}{нет} \\
	\midrule
	Бадзоро &
	\multicolumn{3}{c}{2} \\
	\multicolumn{4}{L{0.93\linewidth}}{\small{\textit{Пояснение: бадзоро --- базовая добавочная стоимость в ханах на каждую выигрышную руку}}} \\
	\midrule
	Фуритен риичи &
	\multicolumn{3}{c}{допустимо} \\
	\midrule
	Закрытый кан после риичи, при условии что интерпретация руки не меняется &
	\multicolumn{3}{c}{допустимо} \\
	\midrule
	Рюиисо требует хотя бы пары зеленых драконов &
	\multicolumn{3}{c}{нет} \\
	\midrule
	Сочетаемость риншана и хайтея &
	\multicolumn{3}{c}{нет} \\
	\midrule
	Формальный темпай (без яку) &
	\multicolumn{3}{c}{допустимо} \\
	\midrule
	Если все тайлы, требуемые для победы, уже лежат в открытых сетах  этого же игрока&
	\multicolumn{3}{c}{нотен} \\
	\midrule
	+2 фу за победу по цумо &
	\multicolumn{3}{c}{да, кроме пин-фу} \\
\end{tabularx}
\section{Заключение}

В данном регламенте мы постарались максимально охватить как все вопросы, регулярно возникающие у начинающих игроков, так и нюансы в процессе проведения турниров и клубных игр. В случае возникновения каких-либо вопросов или уточнений, просим обращаться в рабочую группу по правилам, все вопросы обязательно будут рассмотрены и при необходимости будут внесены правки.

Редакции правил предполагается обновлять раз в год. В случае отсутствия существенных правок, выпуск новой редакции может быть отменен. При проведении турниров просим указывать версию (год выпуска) регламента правил, которые вы собираетесь использовать на турнире. Ожидается, что все судьи турнира ориентируются в регламенте и следят за актуальностью своих знаний. Рабочая группа обязуется обеспечить максимальное публичное распространение новых редакций по мере их выхода.

Спасибо за внимание и ждем вас на турнирах!
